% Generated draft: requires mathematical and source review.
\documentclass[11pt]{article}
\usepackage[margin=25mm]{geometry}
\usepackage{fontspec}
\setmainfont{Latin Modern Roman}
\usepackage{amsmath,amssymb,mathtools,unicode-math}
\setmathfont{Latin Modern Math}
\usepackage{enumitem,needspace,xurl,hyperref}
\setlength{\parindent}{0pt}
\setlength{\parskip}{5pt}
\setlength{\emergencystretch}{3em}
\newcommand{\N}{\mathbb{N}}
\newcommand{\Z}{\mathbb{Z}}
\newcommand{\Q}{\mathbb{Q}}
\newcommand{\R}{\mathbb{R}}
\newcommand{\C}{\mathbb{C}}
\newcommand{\F}{\mathbb{F}}
\newcommand{\A}{\mathbb{A}}
\newcommand{\PP}{\mathbb{P}}
\newcommand{\E}{\mathbb{E}}
\newcommand{\eps}{\varepsilon}
\newcommand{\dd}{\,\mathrm{d}}
\begin{document}
\section*{UnsolvedMath ID: 30000449}
\textbf{Source problem code:} catalogue entry
\subsection{Definitions and mathematical statement}
The source asks whether five-term cohomology sequences for extensions of categories are the low-degree terms of Lyndon--Hochschild--Serre-type spectral sequences, as for group extensions. The categorical extension notion, coefficient system and existing five-term maps are not reproduced.

For comparison, a group extension $1\to K\to G\to Q\to1$ and a suitable $G$-module $M$ give
\[
 E_2^{p,q}=H^p(Q;H^q(K;M))\Longrightarrow H^{p+q}(G;M).
\]
In a small-category framework one may define category cohomology with a covariant abelian-group coefficient functor as right-derived inverse limits; contravariant conventions require the opposite category. For a specified extension $\mathcal C\to\mathcal D$, a categorical analogue would require coefficient systems $\mathcal H^q$ on $\mathcal D$ formed from the appropriate kernel/fiber cohomology, and a convergent first-quadrant spectral sequence
\[
 E_2^{p,q}=H^p(\mathcal D;\mathcal H^q)
                  \Longrightarrow H^{p+q}(\mathcal C;M).
\]
The expected low-degree exact sequence is
\[
 0\longrightarrow E_2^{1,0}\longrightarrow H^1(\mathcal C;M)
 \longrightarrow E_2^{0,1}\xrightarrow{\ d_2\ }E_2^{2,0}
 \longrightarrow H^2(\mathcal C;M).
\]
The requested identification must recover the source's particular five-term groups \emph{and maps}, including edge and transgression maps, not merely produce some unrelated spectral sequence with a formally similar shape.

For categories, specifying an arbitrary functor does not by itself identify a group-like kernel, fiber cohomology action or the proposed extension hypotheses. Those data are essential to the displayed analogue. The record does not fix finite versus general small categories, the extension axioms, coefficient variance or convergence conditions.

The supplied literature assessment checked 21 August 2026 attributes LHS-type spectral sequences to Xu and Webb under an explicit extension hypothesis, while warning that this does not establish the unrestricted categorical generality suggested by the extraction. That hypothesis is not supplied, so the partial theorem is retained with its restriction and not generalized to every category functor.
\subsection{Short English statement}
Do the known five-term sequences for category extensions come from a full LHS-type spectral sequence with the same edge and transgression maps? The extension and coefficient conventions are missing. The supplied assessment reports Xu--Webb results under a specific extension hypothesis, without an unrestricted theorem.
\subsection{Sources}
\begin{itemize}
\item Local UnsolvedMath registry, permanent ID 30000449, exact source code catalogue entry; \emph{Spectral Sequences for Extensions of Categories}. Proposed year 2006, status partially\_solved; corrected\_verified editing only normalizes the theorem-name typography.
\item Supplied report citation: \emph{Representations of Finite Groups (2007)}, Oberwolfach Reports, DOI 10.4171/owr/2006/15, pp. 895--978, \url{https://doi.org/10.4171/owr/2006/15}. Title year differs from DOI/proposed year.
\item Supplied literature assessment checked 2026-08-21: ``OWR report Representations of Finite Groups'', \url{https://doi.org/10.4171/owr/2006/15}; ``Webb low-dimensional cohomology of categories'', \url{https://www-users.cse.umn.edu/~webb/Publications/OWMarch06.pdf}. It attributes the restricted spectral-sequence result to Xu and Webb. The dated OWR-derived curation requests independent verification of time-sensitive status and cited claims; no linked PDF was accessed here.
\end{itemize}
\end{document}
